\section{Introduction}

\label{sec:introduction}
% ══════════════════════════════════════════════════════════════════════════════

Climate change is reshaping the geographic ranges of species at
unprecedented rates \citep{chen2011rapid, pecl2017biodiversity}. As
temperature and precipitation regimes shift, species must either adapt in
situ, migrate to newly suitable habitats, or face local extinction
\citep{urban2015accelerating}. Understanding and predicting these
distributional changes is essential for effective conservation planning,
protected area design, and biodiversity policy.

Species distribution models (SDMs) have become the primary computational
tool for projecting range shifts \citep{elith2009species,
guisan2005predicting}. Traditional approaches, particularly MaxEnt
\citep{phillips2006maximum} and generalized additive models
\citep{hastie1990generalized}, establish correlative relationships between
species occurrence records and environmental covariates. While widely
adopted, these models suffer from several well-documented limitations:

\begin{enumerate}[label=(\roman*)]
  \item \textbf{Linearity assumptions}: Most correlative SDMs struggle to
    capture the complex, non-linear interactions between environmental
    variables that determine habitat suitability
    \citep{elith2006novel}.
  \item \textbf{Spatial independence}: Standard SDMs treat occurrence
    records as independent samples, ignoring spatial autocorrelation and
    dispersal constraints \citep{dormann2007methods}.
  \item \textbf{Static landscapes}: Projections typically assume
    instantaneous equilibrium with climate, ignoring transient dynamics,
    dispersal limitation, and habitat fragmentation
    \citep{zurell2016benchmarking}.
  \item \textbf{Connectivity neglect}: The physical connectivity of
    landscapes---whether corridors exist for species to track shifting
    climate envelopes---is rarely incorporated into distributional
    predictions \citep{hodgson2009climate}.
\end{enumerate}

Recent advances in deep learning offer pathways to address these
shortcomings. Neural network-based SDMs can model arbitrary non-linear
relationships \citep{botella2018deep, chen2020deep}, and graph neural
networks (GNNs) provide a natural framework for spatially structured
ecological data \citep{jetz2019essential}.

\subsection{Contributions}

This paper introduces \textbf{HabitatNet}, a graph neural network
framework for species distribution modeling that addresses the above
limitations. Our contributions are:

\begin{enumerate}
  \item A \textbf{spatial graph representation} of landscapes where
    nodes encode local environmental features at 1\,km resolution and
    edges represent dispersal potential weighted by terrain resistance.

  \item A \textbf{multi-species GNN architecture} that shares learned
    environmental representations across taxa while learning
    species-specific habitat preferences through attention-weighted
    readout layers.

  \item Integration with \textbf{CMIP6 climate projections} under
    multiple SSP scenarios (SSP1-2.6, SSP2-4.5, SSP3-7.0, SSP5-8.5)
    for temporal forecasting through 2100.

  \item A \textbf{corridor identification algorithm} built on learned
    graph embeddings that identifies critical connectivity bottlenecks
    and proposes optimal corridor placements.

  \item Comprehensive evaluation on \textbf{847 terrestrial vertebrate
    species} across six biogeographic realms, with temporal validation
    against independent resurvey data.
\end{enumerate}

\subsection{Relation to Zach Kelling}

Zach Kelling operates as an open AI research network dedicated to
decentralized science (DeSci) and decentralized AI (DeAI). This work
advances the Foundation's mission by: (a) developing open-source AI
tools for biodiversity conservation, (b) establishing reproducible
benchmarks through the Zoo Open Science Network, and (c) enabling
community-driven conservation planning through transparent, verifiable
models. All artifacts are governed through Zoo Improvement Proposals
(ZIPs) at \texttt{zips.zoo.ngo}.


% ══════════════════════════════════════════════════════════════════════════════
